% Section partagée par les deux rapports M4.3Live : le modèle et les
% notations, pour que chaque document se lise sans accès au code ni aux
% rapports des lignées antérieures.
\section{Le modèle, et les notations employées}
\label{sec:modele}

\subsection{Ce que simule M4.3Live}

Le modèle décrit une population d'entités économiques abstraites qui
échangent une ressource unique, le \emph{joule}, assimilée ici à un capital.
Il n'y a ni monnaie, ni banque centrale, ni prix : seulement des capitaux,
une fonction de production, et des prêts bilatéraux perpétuels.

Une entité $i$ est entièrement décrite par trois nombres : son capital
$K_i \ge 0$ et sa \emph{technologie} $(A_i, \gamma_i)$, qui fixe la quantité
qu'elle produit à chaque pas de temps,
\begin{equation}
\text{production de } i \;=\; A_i\,K_i^{\gamma_i},
\qquad A_i > 0,\quad 0 < \gamma_i < 1 .
\label{eq:production}
\end{equation}
L'exposant $\gamma_i < 1$ rend les rendements \emph{marginaux} décroissants :
un joule supplémentaire rapporte d'autant moins que le capital est déjà
grand ($\partial^2/\partial K^2 < 0$). La production, elle, continue de
croître avec le capital --- doubler le capital augmente la production, mais
de moins du double. C'est cette concavité qui crée un intérêt à déplacer du
capital d'une entité riche vers une entité pauvre --- le fondement même du
marché du crédit du modèle.

\subsection{Un pas de temps}

Chaque pas déroule huit phases, dans cet ordre invariable :

\begin{enumerate}
\item \textbf{Interventions} --- les changements de paramètres demandés en
  direct sont appliqués (propre à M4.3Live ; voir §\ref{sec:portees} pour
  leurs portées).
\item \textbf{Naissances} --- un nombre $\mathrm{Poisson}(\lambda)$ d'entités
  neuves apparaissent, chacune dotée du capital de naissance $K_0$.
\item \textbf{Choc} --- chaque capital est multiplié par $e^{\xi}$ avec
  $\xi \sim \mathcal{N}(-\sigma^2/2,\ \sigma^2)$, un choc multiplicatif
  d'espérance 1 (il ne crée ni ne détruit de richesse en moyenne).
  \emph{Ordre de grandeur.} L'écart typique se lit sur
  $\mathbb{E}|\xi| = \sigma\sqrt{2/\pi}$ à la correction de moyenne près
  ($-\sigma^2/2$, négligeable ici) : à $\sigma = 0{,}01$, la valeur employée
  dans tout ce programme, $\mathbb{E}|\xi| \approx 8{,}0\times10^{-3}$, donc
  un pas multiplie un capital par $1 \pm 0{,}8\,\%$ en moyenne, et par
  $1 \pm 2{,}0\,\%$ dans 5\,\% des cas. Le choc est un bruit de fond, pas un
  moteur : sur les 2000 pas d'un amorçage, sa contribution \emph{nette}
  cumulée (\code{shock\_gain}) vaut $-0{,}41\,\%$ du capital agrégé final,
  et aucun pas isolé ne déplace \code{K\_tot} de plus de $0{,}16\,\%$.
\item \textbf{Production} --- chaque entité produit selon
  \eqref{eq:production} et \emph{ajoute intégralement} sa production à son
  propre capital.
\item \textbf{Service des intérêts} --- chaque emprunteuse doit
  $\mathrm{due} = \sum_k q_k r_k$, somme sur ses contrats du principal $q_k$
  fois le taux $r_k$. Si son capital ne suffit pas, elle paie ce qu'elle
  peut, tombe à zéro et est marquée en \emph{défaut de liquidité}.
\item \textbf{Dépréciation} --- chaque capital est multiplié par $1-\delta$.
\item \textbf{Marché du crédit} --- $\lfloor \eta(N) \rfloor$ tirages de
  paires, où $N$ est le nombre d'entités éligibles et
  $\eta_{\rho,\beta}(N) = \rho\,N_{\text{ref}}(N/N_{\text{ref}})^{\beta}$
  (la baseline $\rho=1$, $\beta=1$ donne simplement $\eta(N)=N$). Dans chaque
  paire, la plus riche est la \emph{prêteuse}, la plus pauvre
  l'\emph{emprunteuse} ; le prêt ne va jamais dans l'autre sens. Un contrat
  transfère un principal $q$ et fixe un taux $r$, tous deux gelés à vie.
\item \textbf{Faillites en cascade} --- toute entité dont la valeur nette
  $K + \text{créances} - \text{dettes}$ devient négative, ou qui est en
  défaut de liquidité, meurt : ses dettes sont annulées (perte sèche pour
  ses prêteuses, qui peuvent tomber à leur tour) et son capital est détruit.
  On itère jusqu'à point fixe.
\end{enumerate}

Un identifiant d'entité n'est jamais réutilisé ; une entité morte le reste.

\subsection{Notations}

\begin{center}\footnotesize
\begin{tabular}{@{}llp{8.6cm}@{}}
\toprule
\textbf{Symbole} & \textbf{Nom} & \textbf{Définition} \\
\midrule
$K_i$ & capital & Ressource détenue par l'entité $i$. \\
$A_i$ & coefficient d'extraction & Facteur multiplicatif de la production
\eqref{eq:production}. \textbf{Levier primaire} de ce programme. \\
$\gamma_i$ & exposant d'extraction & Exposant de la production : il fixe le
rendement d'échelle. C'est lui qui rend les rendements \emph{marginaux}
décroissants ($0<\gamma<1$), et non le niveau de la production, qui croît
toujours avec le capital. Levier secondaire de ce programme. \\
$(A_i,\gamma_i)$ & technologie & Couple fixé à la naissance, jamais modifié
sauf par une intervention. Reçoit un identifiant entier interne. \\
$\lambda$ & taux de naissance & Espérance du nombre d'entités créées par pas. \\
$K_0$ & capital de naissance & Capital dont une entité neuve est dotée. \\
$\delta$ & dépréciation & Fraction du capital perdue à chaque pas. \\
$\sigma$ & amplitude du choc & Écart-type du choc multiplicatif log-normal. \\
$N$ & taille du bassin & Nombre d'entités éligibles au marché à ce pas. \\
$x,\ y$ & capitaux de la paire & $x = K_b$ (emprunteuse), $y = K_\ell$
(prêteuse), avec $y > x$ par la règle de marché. \\
$C$ & somme conservée & $C = x+y$ ; le transfert la laisse inchangée. \\
$q$ & principal & Montant transféré par le contrat. \\
$r$ & taux & Service versé par pas, perpétuel : l'emprunteuse doit $q\,r$ à
chaque pas, indéfiniment. \\
$\delta^{*}$ & transfert optimal & Principal qui maximise la production
jointe de la paire (§\ref{sec:modele}, éq.~\ref{eq:transfert}). \\
$\lambda^{*}$ & part optimale & Part de $C$ revenant à l'emprunteuse à
l'optimum ; $h(C) = C\lambda^{*}$ est son capital optimal. \\
$\Kaut$ & échelle autarcique & $[A(1-\delta)/\delta]^{1/(1-\gamma)}$ : capital
auquel production et dépréciation s'équilibrent sans marché. Repère
d'échelle du système. \\
\code{prod\_tot} & production agrégée & Somme des productions du pas.
\textbf{Variable d'intérêt de ce programme.} \\
\code{K\_tot}, \code{pop} & agrégats & Capital total et effectif vivant. \\
\code{deaths} & morts & Entités mortes au pas (racines + victimes de cascade). \\
\code{defaults} & défauts & Entités en défaut de liquidité au pas. \\
\code{roots\_insolvency} & racines d'insolvabilité & Entités dont la valeur
nette est passée sous zéro par leurs propres pertes (à distinguer des
victimes de cascade). \\
\code{loan\_volume} & volume de prêt & Somme des principaux transférés au pas. \\
\code{mkt\_rounds} & tirages & Nombre de paires tirées au pas. \\
\code{mkt\_blocked\_dir} & paires refusées & Paires où $\delta^{*} \le 0$ :
l'optimum voudrait faire circuler le capital du pauvre vers le riche, ce que
le sens du prêt interdit. \textbf{Compteur propre à M4.3Live.} \\
\bottomrule
\end{tabular}
\end{center}

\subsection{L'institution de principal de M4.3Live}

C'est le seul point où M4.3Live diffère de son prédécesseur M4.3. Là où M4.3
transférait la moitié de l'écart de capital, $q = (y-x)/2$, M4.3Live
transfère le montant qui maximise la production \emph{jointe} de la paire
immédiatement après l'échange :
\begin{equation}
\delta^{*} \;=\; \arg\max_{-x \le \delta \le y}\;
A_b\,(x+\delta)^{\gamma_b} + A_\ell\,(y-\delta)^{\gamma_\ell}
\;=\; h(C) - x,
\qquad h(C) = C\,\lambda^{*}.
\label{eq:transfert}
\end{equation}
Trois conséquences importent pour lire les deux rapports :

\begin{enumerate}
\item Quand les deux entités partagent la \emph{même} technologie,
  $\lambda^{*} = 1/2$ et $\delta^{*} = (y-x)/2$ exactement : la baseline
  homogène de M4.3Live \emph{est} celle de M4.3.
\item Dès que les technologies diffèrent, $\lambda^{*} \ne 1/2$ et le partage
  est pondéré en faveur de celle dont le \emph{rendement marginal est le plus
  élevé au capital considéré}. Il n'existe pas de « meilleure »
  technologie dans l'absolu : à exposants inégaux, l'ordre des rendements
  marginaux $\gamma A K^{\gamma-1}$ s'inverse avec le capital. Un $\gamma$
  élevé l'emporte aux grands capitaux, un $A$ élevé aux petits, et il existe
  toujours un capital de croisement. C'est pourquoi $\lambda^{*}$ dépend de
  $C$ et non seulement des paramètres.
\item $\delta^{*}$ peut être \emph{négatif}. Comme $\delta^{*} = C\lambda^{*} - x$,
  on a exactement
  \[ \delta^{*} \le 0 \iff \frac{x}{C} \ge \lambda^{*} , \]
  c'est-à-dire : dès que la prêteuse est celle dont le rendement marginal est
  le plus élevé, l'optimum voudrait lui envoyer du capital, et la paire ne
  traite pas. C'est le compteur \code{mkt\_blocked\_dir}, et c'est un canal à
  part entière, pas un détail numérique.
\end{enumerate}

\paragraph{Surplus coopératif.} Le \textbf{surplus coopératif} d'une paire
est le gain de production \emph{jointe} que l'échange produit, à capital
total inchangé :
\begin{equation}
S(x,y) \;=\;
\underbrace{A_b (x+q)^{\gamma_b} + A_\ell (y-q)^{\gamma_\ell}}
          _{\text{après le transfert}}
\;-\;
\underbrace{A_b\,x^{\gamma_b} + A_\ell\,y^{\gamma_\ell}}_{\text{avant}}
\;\ge\; 0 ,
\label{eq:surplus}
\end{equation}
où $q$ est le principal effectivement transféré --- égal à $\delta^{*}$ sauf
quand le plafond mord. La positivité est immédiate quand $q = \delta^{*}$ :
$\delta^{*}$ maximise le premier terme et $\delta = 0$ est admissible ; sous
plafond elle reste vraie car $q$ est entre $0$ et $\delta^{*}$ et le premier
terme est concave en $\delta$. $S$ est nul exactement quand la paire est
déjà à l'optimum ($x/C = \lambda^{*}$), et les paires refusées ne
contribuent pas. La colonne \code{mkt\_surplus} des séries est la somme de
$S$ sur les paires traitées au pas. C'est un \emph{flux de
production par pas}, homogène à \code{prod\_tot}, et non un stock de
capital : le marché ne crée aucun joule, il déplace du capital vers l'endroit
où il produit davantage.

\subsection{Les trois portées d'intervention}
\label{sec:portees}

Une intervention change un paramètre \emph{pendant} qu'une simulation
tourne. Sa \textbf{portée} dit qui elle touche :

\begin{description}
\item[\code{all} (toutes)] --- rétroactif \emph{et} prospectif : toutes les
  entités vivantes reçoivent la nouvelle valeur, et la valeur par défaut à la
  naissance change aussi.
\item[\code{new} (nouvelles)] --- \emph{vintage} technologique : seule la
  valeur par défaut à la naissance change. Les entités déjà vivantes gardent
  la leur pour toujours ; la population se convertit par renouvellement.
\item[\code{fraction} (une fraction $\varphi$)] --- une fraction $\varphi$
  des entités vivantes est tirée au sort une fois, et elle seule est
  modifiée. \textbf{La valeur par défaut à la naissance n'est pas touchée} :
  aucune naissance ne vient reconstituer la cohorte traitée.
\end{description}

Certaines combinaisons n'ont pas de sens et sont refusées explicitement :
$K_0$ n'existe qu'au moment de la naissance, donc \code{all} et \code{new} y
sont le même objet et \code{fraction} n'a pas d'objet ; $\lambda$, $\delta$,
$\sigma$, $\rho$ sont des paramètres de population et non des attributs
d'entité, donc \code{new} y est un alias explicite de \code{all}.
